\section{Related Work}

\subsection{Tool-Using LLM Agents}

Recent progress in large language models (LLMs) has enabled agents to interact with external tools such as search engines, APIs, calculators, and execution environments. A representative line of work studies how language models can interleave reasoning and action during inference. ReAct is a canonical example, showing that alternating between reasoning traces and external actions can improve performance on knowledge-intensive and decision-making tasks~\cite{yao2023react}. Toolformer further demonstrates that language models can learn when and how to invoke tools in a more structured way~\cite{schick2023toolformer}. Gorilla focuses on API-centric tool use and emphasizes the importance of accurate tool selection under realistic interfaces~\cite{patil2023gorilla}. More recent benchmarks and frameworks such as ToolBench further highlight tool-use planning, tool selection, and interaction quality as central concerns for practical LLM agents~\cite{qin2023toolbench}.

These works establish the feasibility and usefulness of tool-using agents, but most of them primarily focus on whether the agent can successfully use tools, rather than how to explicitly control the cost of multi-step behavior. In contrast, our work studies \emph{agent orchestration} as an explicit policy layer that regulates when additional retrieval, tool use, verification, or stopping is worthwhile.

\subsection{Reasoning, Planning, and Self-Improvement}

A second line of work studies how LLMs can reason more effectively through intermediate deliberation or self-improvement. Chain-of-Thought prompting shows that explicit intermediate reasoning can significantly improve complex problem solving~\cite{wei2022chain}. Self-consistency further improves reasoning by sampling multiple chains and aggregating their answers~\cite{wang2023selfconsistency}. Tree of Thoughts generalizes linear reasoning into a more structured search over intermediate thoughts~\cite{yao2023tree}. Program-of-Thoughts prompting instead leverages executable programs to support more reliable reasoning~\cite{chen2023program}. Other works focus on iterative refinement and self-correction, including Self-Refine~\cite{madaan2023selfrefine}, Reflexion~\cite{shinn2023reflexion}, and LATS, which combines language models with tree search over action trajectories~\cite{zhou2024lats}.

While these methods improve reasoning quality, they also tend to increase the number of intermediate steps and the amount of accumulated context. As a result, they do not directly answer the systems-level question of when more reasoning is actually justified. Our work is complementary: rather than proposing a stronger reasoning paradigm, we focus on explicit action selection and stopping decisions under quality--cost trade-offs.

\subsection{Agent Systems, Orchestration, and Multi-Agent Frameworks}

A growing body of work studies LLM agents as systems rather than isolated prompting techniques. AutoGen investigates multi-agent conversation as a mechanism for task decomposition and collaboration~\cite{wu2023autogen}. MetaGPT explores role-based multi-agent software workflows~\cite{hong2024metagpt}. ChatDev similarly models software development as multi-agent communication among role-specialized agents~\cite{qian2024chatdev}. In a different direction, Voyager studies lifelong embodied agents that accumulate skills through iterative exploration~\cite{wang2023voyager}. These works reflect a broader shift from single-shot prompting to agentic systems that involve memory, interaction, coordination, and environment feedback.

However, many existing agent systems still rely heavily on fixed workflows, hand-designed loops, or prompt-driven control. They often leave the orchestration logic implicit, making it difficult to analyze how execution cost arises and how different control signals affect the resulting trajectories. Our work differs in that it makes the orchestration layer explicit and lightweight: instead of proposing a full multi-agent framework or a new interactive environment, we focus on a utility-guided decision mechanism for single-agent tool use.

\subsection{Efficient and Cost-Aware Inference}

Efficiency has become an increasingly important issue in modern LLM systems. Prior work has explored efficient transformer architectures~\cite{tay2022efficient}, adaptive computation mechanisms~\cite{graves2016adaptive}, and early exiting strategies for reducing inference cost~\cite{elbayad2020depth}. Other works study selective prediction, confidence-based routing, or budget-aware computation for neural models~\cite{schuster2022confident,xin2020deebert}. More broadly, recent LLM deployment research has emphasized practical constraints such as token budgets, latency, and context growth.

Compared with this line of research, our work does not propose a new model architecture or a learned adaptive-computation module. Instead, we study efficiency at the orchestration level: the agent explicitly evaluates whether another step is likely to be beneficial, and balances estimated gain against cost-related signals. This makes our approach closer to a lightweight policy layer for tool-using LLM agents, with a focus on analyzing quality--cost trade-offs rather than maximizing task performance alone.

\subsection{Agent Benchmarks and Evaluation}

Recent work has also emphasized the need for realistic evaluation settings for LLM agents. WebShop provides an interactive environment for evaluating multi-step decision making in a web-based shopping setting~\cite{yao2022webshop}. ALFWorld evaluates embodied instruction-following agents in text-based household environments~\cite{shridhar2021alfworld}. ToolBench introduces a large-scale benchmark specifically designed for tool-use agents~\cite{qin2023toolbench}. Other works have pointed out that agent evaluation should consider not only final answer quality, but also trajectory quality, action efficiency, and robustness under realistic execution settings.

Our evaluation perspective is aligned with this broader trend. Rather than measuring only final task quality, we also examine token consumption, wall-clock time, redundancy, and workflow fairness. In this sense, our work contributes to the emerging view that agent quality should be studied together with the computational and behavioral cost of obtaining that quality.
Overall, unlike prior work that primarily emphasizes stronger reasoning, broader tool access, or richer agent frameworks, our work focuses on explicit orchestration as a lightweight decision layer for analyzing and controlling the quality--cost trade-off of tool-using LLM agents.